\chapter{对抗安全}
\label{chap:trust-security}

邮件分类器可能因新月份的语言变化而误判，也可能因为有人专门改写邮件来逃避检测而失效。后者需要主动寻找薄弱点，并选择有利于攻击目标的输入或操作。若分类器进一步接入自动归档、检索和发送工具，攻击造成的后果就可能从判断错误扩大到信息泄露或越权执行。

本章讨论机器学习系统怎样抵御这种有目的的操纵。攻击的定义、例子与损害用于确定保护目标；脆弱性评估检查输入、训练过程、模型内部、查询接口和执行链中的薄弱环节；防御则在相应位置减少攻击机会、限制模型受影响的程度，并控制已经发生的损害。分布偏移（第~\ref{sec:distribution-shift}节）描述训练分布与评价分布的差异，算法稳定性（第~\ref{sec:algorithmic-stability}节）考察训练样本变化对输出或损失的影响，隐私泄露（第~\ref{sec:privacy-leakage-channels}节）关注系统暴露了什么信息；这些现象即使没有攻击者也可能发生。对抗安全进一步要求有人在其知识、权限和资源范围内有目的地选择变化，并追踪这种操纵怎样经过系统响应和下游动作形成损害。本章据此检验整个系统能否守住规定的安全边界。

\section{对抗安全的基本概念}
\label{sec:security-threats}

\subsection{对抗攻击及其典型表现}
\label{sec:security-attack-definition}

\term{对抗攻击}（adversarial attack）是攻击者有目的地选择或操纵机器学习系统能够接触的输入、训练材料、模型组件或交互过程，使系统违背预期任务或安全要求的行为。它既包括让分类器作出错误判断，也包括获取不应公开的信息、不当占用服务资源或诱导越权操作。本章采用这一较宽的含义；狭义的输入扰动攻击只是其中一种。

攻击与普通错误的区别在于有目的的选择和利用。系统偶然误判一封邮件，只能说明存在错误；有人反复试探并选择能够逃过过滤的写法，才构成这里讨论的攻击过程。反过来，攻击已经发生也不意味着攻击已经成功：输入可以被拒绝，错误建议可以被执行器拦下，实际损害还要另行判定。

攻击入口是攻击者可以施加影响的位置。下面沿同一个邮件系统说明几类典型表现；这些是帮助区分机制的教学情境，不是已发生事件的记录。

输入攻击改变模型看到的内容。\term{对抗样本}（adversarial example）在规定修改范围内诱导不期望的预测，范围可以是数值扰动，也可以是任务允许的文本或物理变化。修改是否保留原任务标签需要核验，否则测试可能混入任务本身已经改变的输入。例如，攻击者改写钓鱼邮件的措辞，使分类器将其判为正常邮件，而邮件的欺骗目的仍未改变。

训练阶段的\term{数据投毒}（data poisoning）通过修改训练数据改变学到的模型；\term{后门}（backdoor）使模型在特定触发条件下表现异常，同时可能保留正常输入上的性能。干净测试集准确率因而不足以排除后门。联邦或协作训练中，恶意参与者提交的更新也属于训练入口，需要明确其比例和可控内容。例如，污染反馈数据使某类钓鱼邮件被当作正常样本学习，或让带特定标记的邮件在部署后被放行。后门既可以通过数据投毒植入，也可以通过受污染的模型参数进入系统，因此投毒描述操作过程，后门描述被植入的异常行为。

模型服务还暴露查询与输出。模型窃取关注攻击者能否复制有价值的功能或信息，其目标不同于第~\ref{sec:privacy-concepts}节的个人数据泄露。两者可能使用相似查询接口，却需要不同的成功定义和防御代价评价。例如，攻击者在预算内反复查询分类器，再训练一个具有近似判断能力的替代模型；私人邮件内容被恢复则属于另一种信息损害。

对于具有工具和检索能力的生成系统，\term{提示注入}（prompt injection）把攻击指令放进模型要读取的内容中，试图改变原任务；间接注入将这些内容置于第三方文档、网页或工具返回值。\term{越狱}（jailbreak）则试图使模型突破既定行为限制。二者可能结合，但来源、目标和系统权限仍需分别说明。例如，助手读取一封邮件以生成摘要，正文却夹带要求归档其他邮件的指令。它利用的是系统对内容角色和操作权限的处理，而不要求对输入施加肉眼不可见的扰动。

\subsection{攻击影响与损害路径}
\label{sec:security-damage-paths}

损害目标需要与正常任务区分。逃过分类、诱导指定错误、提取训练信息、复制模型功能和执行未经允许的工具动作，是不同目标。若只把所有失败合成一个“攻击成功率”，就无法知道主要损害来自哪里，也难以选择防御位置。

损害分析还要区分单个样本与整个服务。攻击者可能希望让某封钓鱼邮件通过，也可能希望大量正常邮件被隔离，从而使用户关闭过滤功能。前者破坏决策的完整性，后者损害服务的可用性；若攻击使助手泄露其他邮件，则涉及机密性。把三种损害分开，才能解释为什么提高邮件获得放行所需的可信分数阈值，可能抑制第一种攻击，却加剧第二种损失。

例如，用户授权助手归档指定类别的邮件，某封邮件正文却声称“应立即归档全部邮件”。这段文字可以是待分析数据，却没有因此获得操作其他邮件的权限。攻击成功依赖系统把内容提升为行动依据，而不只依赖模型生成了某句话。

一次攻击通常沿“可控入口—系统响应—下游动作—损害”传播。输入修改是否令分类器翻转、训练污染是否改变后来学到的规则、工具返回值是否诱导越权动作，是不同的中间事件。评估应沿实际路径观察，不能仅凭中间事件替代损害结果。例如，模型提出越权操作说明行为受到操纵，但只有执行端实际允许该操作，才产生相应行动损害。

同一目标也可能有多条路径。让一封钓鱼邮件通过，既可以改写这封邮件，也可以提前污染训练数据；两者都损害决策完整性，却需要不同权限、预算和测试过程。这说明攻击入口与损害目标是交叉关系，评估既要覆盖不同入口，也要使用与各自损害对应的判据。

可用性攻击还可能主要消耗资源，而不改变答案。例如，在受控测试中让请求持续占用计算、使任务陷入重复调用，可以观察队列等待、正常请求完成率和资源开销是否恶化。这类试验应报告单位请求及整个攻击过程的成本，并与同负载的正常使用比较，避免把普通容量不足误认成特定攻击机制。

\subsection{攻击能力与威胁模型}
\label{sec:security-capabilities}

\term{威胁模型}（threat model）规定攻击者想达到什么结果、能观察和改变什么，以及受到什么资源限制。这些条件共同限定安全结论：攻击入口说明通过哪里施加影响，知识与权限说明可以如何操作，损害目标说明什么结果才算成功。同一次攻击可以同时具有特定入口、知识条件和损害目标，因此不能把三个维度的取值列成互斥类别。

在给定威胁模型下，安全要求可以表述为：允许的攻击过程不应使系统产生规定的损害，或者其损害风险应低于容许水平。“允许”指测试给予攻击者的能力范围，并非给予其现实行动许可。例如，测试可以允许攻击者改写邮件内容，却不允许直接修改收件箱数据库。前者需要检查模型与执行机制如何响应，后者若被错误地纳入测试，便绕过了原本希望评价的防护环节。

白盒攻击可以访问模型参数或梯度，黑盒攻击可能只得到查询结果；实际系统还可能向攻击者暴露部分结构、版本和错误信息。知道参数与能够修改参数也应区分，前者可以帮助寻找输入，后者则可能直接改变模型。

资源限制包括查询次数、修改幅度、能够污染的数据量和可控制的参与者数量。较强权限下找到的攻击不一定能在较弱权限下复现，但它可以揭示机制缺口。反之，在很小预算下未找到攻击，也不能说明扩大预算后仍安全。

白盒与黑盒描述的是知识条件，不是攻击的善恶或现实程度。设计者用完整模型进行白盒测试，是为了排除“防御只依赖攻击者暂时不了解系统”的可能。公开服务还需要黑盒测试，因为只能观察类别、能够观察概率和能够访问内部表示，提供的信息量并不相同。评估中的权限必须真实用于搜索，不能给攻击者白盒访问，却只运行一个没有充分利用梯度的简单方法。

\section{脆弱性的评估}
\label{sec:security-evaluation}

脆弱性评估要判断：在给定能力和预算下，攻击者能通过哪些路径造成何种损害。输入测试、模型内部分析与系统审查可以互相补充，并不是只能选择其中一种。以下先统一评估维度、指标与试验程序，再按被测对象展开五类检查：固定模型而改变输入，改变材料并重新运行训练，读取已有模型的内部机制，查询已发布接口，以及运行带权限和状态的完整任务。

这些对象可以属于同一条攻击链。例如，后门可能由训练材料进入参数，再由特定输入触发，并通过工具执行造成损失。分开评估有助于定位薄弱环节，最终结论仍需回到完整系统。白盒搜索、黑盒试探、组件核查和形式化验证是取得证据的方法；选哪一种，取决于相应对象能提供什么信息以及要判断什么性质。

\subsection{评估维度与判定标准}
\label{sec:security-risk-assessment}

评估至少需要区分被测对象、攻击条件与结果指标。输入、模型内部和工具链说明检查哪里；白盒或黑盒、查询次数和污染比例说明在什么能力下检查；攻击成功率、损害和正常任务表现说明检查得到了什么。它们是可以组合的维度，例如同一个输入入口既可以做白盒测试，也可以做黑盒测试，并分别测量不同预算下的损害。

结果指标应回答四类问题：有多少任务被攻破，攻破需要多少修改或查询，成功后影响多大，以及防护使多少正常任务受损。评估覆盖范围则说明这些数字适用于哪些攻击、数据与系统配置。下文的攻击成功率刻画第一项，预算曲线与污染强度刻画第二项，实际状态变化和恢复成本刻画第三项，干净任务效用、误拒率与延迟刻画第四项。它们不能合成一个未经说明的“安全分数”。

\paragraph{可重复的评估程序}

基于攻击试验的评估通常包括准备、搜索、执行核验和结果比较四个相互依赖的步骤。准备阶段规定攻击者实际能得到的信息；搜索阶段构造候选攻击；执行核验确认候选满足权限和预算，并观察真实系统；结果比较则把目标达成、损害和正常任务代价分开。后面各类攻击试验共享这套程序，但需要取得的材料、可使用的搜索方法和核验对象有所不同。内部核查和形式化验证还可以直接检查计算结构或约束，不一定经过候选攻击搜索；它们应按各自的证据范围报告结论。

\begin{enumerate}
 \item 固定被测模型或训练算法的版本、预处理、权限配置及正常任务，保留未攻击时的输出或运行结果。选定攻击入口，并把可修改对象、可见反馈、预算和成功事件写成可检查条件。
 \item 只向测试攻击者提供约定的信息，让其在预算内生成候选。记录每次查询与反馈、采用的搜索设置，以及最终提交的对象；改变预算后应作为另一个设置单独报告。
 \item 把有效候选交给完整系统执行。既核验修改是否仍属于允许范围，也检查最终输出或状态是否达到攻击目标。对随机系统重复完整过程，保留失败、中止和被拒绝的运行。
 \item 按事先定义的试验单位统计成功与损失，并与正常任务对照。改动防御时，先复用原攻击检查已知问题，再开放同预算的新搜索，评价攻击者适应防御后的结果。
\end{enumerate}

例如，若一轮试验允许修改一封邮件、查询二十次，那么统计单位是这次二十次预算内的完整搜索，而不是把二十次调用当成二十个独立任务。若目标是越权归档，只有授权集合外的邮件实际发生状态变化才算该目标成功；生成一句越权建议可以另记为模型行为偏离。

设$n$次完整试验的成功指示为$s_i\in\{0,1\}$，则经验攻击成功率为$\widehat{\mathrm{ASR}}=n^{-1}\sum_i s_i$。这一定义要求先确定哪些试验进入分母。若考察原本能够正确处理的任务，应先确定这组任务，再统计攻击后失败比例；否则正常失败可能被误计为攻击效果。定向攻击还需要检查是否达到指定结果，而不只是产生任意错误。

\paragraph{攻击结果、损害与证据范围}

报告成功率时还应给出试验数和估计误差。独立同分布的完整试验可以使用二项比例区间；若同一任务进行了多次相关搜索，则应保留任务这一统计单位，不能靠重复查询制造更大的独立样本量。搜索、挑选攻击和最终比较最好使用职责明确的任务集合，避免只对攻击构造时已经成功的样本报告结果。

成功次数也不等于损失。一次误分类与一次大范围数据泄露可能具有不同严重程度。令随机
元素$\Xi$汇总一次完整试验中的任务实例、环境状态和系统内部随机性；若攻击策略自身
随机化，其随机数也并入$\Xi$。允许集合$\mathcal A$中的每个$a$都是在抽取$\Xi$前
固定的策略规则，只能读取威胁模型规定的信息。规则可以在执行过程中依据获准观察的输入
和反馈历史选择后续动作，因此“自适应”必须通过策略的信息接口表达，而不是在看到完整
试验结果后再逐结果挑选动作。记相应损失为$L(a,\Xi)$，则
\begin{equation}
 R_{\mathrm{adv}}=\sup_{a\in\mathcal A}
 \mathbb E_{\Xi}[L(a,\Xi)]
 \label{eq:security-adversarial-risk}
\end{equation}
表达规定攻击能力和信息条件下的最不利风险。若攻击者只能预先选择一个固定动作，
$\mathcal A$就应限制为常值策略；若能够观察当前输入或此前查询反馈，则把这些信息
加入策略接口。两种设定一般给出不同风险，不能在未修改$\mathcal A$时把上确界移入
期望。实际测试通常只覆盖$\mathcal A$的一部分，因此得到的是已发现风险的证据，而不是未知攻击风险的精确估计。

应同时报告影响范围、恢复成本和正常任务效用。拒绝所有输入可以阻止许多攻击，也会使系统失去服务能力。适当比较是在相同任务和资源条件下，检查防御减少了哪些损害、增加了多少误拒和延迟。

重复尝试进一步影响比较口径。独立采样时，多次查询至少成功一次的概率会随预算增加，相关计算见第~\ref{sec:privacy-frontiers}节；安全测试还需考虑更强的情形：攻击者根据失败反馈更新下一次输入，尝试之间不再独立且成功概率可能改变。此时应以完整攻击过程为试验单位，报告规定预算内是否达到目标。只报告一次调用的失败比例，无法代表允许反复试探的服务条件。

加入防御后，攻击者可能改变策略。例如，关键词过滤会使攻击者改写表达，监控器会成为新的操纵对象。因此应向测试攻击者提供相应的防御知识或试探机会，再比较结果。防御前的攻击集可以作为回归测试，却不能独立代表防御后的安全边界。

形式化认证提供另一类证据：在规定输入集合或程序条件下证明性质成立。其优势是覆盖整个被规定的集合，限制则在于模型、输入空间和权限假设是否覆盖实际系统。第~\ref{sec:robustness-certification}节的局部变化界只认证相应数值扰动，不能自动覆盖数据投毒或工具越权。

\subsection{推理输入的评估}
\label{sec:security-prediction-test}

普通测试询问一个输入是否预测正确，对抗测试则追问：在允许的修改范围内，是否存在一个仍应得到相同答案、却能使模型出错的输入？Szegedy等通过输入优化展示了神经网络的这种局部脆弱性，并观察到部分扰动能够迁移到不同网络。由此暴露的问题是，独立测试集上的正确分类，并没有穷尽每个样本周围的决策行为\scite{szegedy2014intriguing}

这一测试需要原始输入及其真实标签、完整预测接口和允许修改集合。白盒设置还提供计算梯度所需的模型信息；黑盒设置则限定返回类别或概率。测试者固定模型，只改变输入，候选必须同时满足修改预算、输入合法性和任务标签保持条件，然后在真实接口上重新预测。

\paragraph{利用输入与输出关系搜索反例}

这种脆弱性不必完全归因于复杂的非线性。考虑线性打分$s(\bm{x})=\bm{w}^{\top}\bm{x}+b$，其中$\bm{x},\bm{w}\in\mathbb R^d$。若每个坐标允许改变不超过$\epsilon$，取$\delta_j=\epsilon\operatorname{sign}(w_j)$，打分变化就是$\epsilon\|\bm{w}\|_1$。即使每个坐标的变化很小，多个方向一致的贡献也可能跨越分类边界。这个计算没有证明所有模型都会失败，但说明“单个特征改动小”与“总决策变化小”并不等价。

若要把原本为正的分数压到负侧，可以选择相反的扰动方向。例如，构造二维分类分数$s(\bm{x})=x_1-x_2+0.05$，以正分数预测类别1。输入$(0.5,0.5)$的分数为0.05；在每个坐标允许变化0.1的条件下，候选$(0.4,0.6)$的分数变为$-0.15$，预测翻转。若任务定义确认这个邻域标签仍为1，便找到一个有效反例。若正确标签也随输入改变，翻转就不构成这里要测量的攻击成功。

Goodfellow等据此用局部线性近似构造\term{快速梯度符号法}（fast gradient sign method, FGSM）。设输入损失梯度为$\bm{g}=\nabla_{\bm{x}}\ell(f(\bm{x}),y)$，在无穷范数约束下，
\begin{equation}
 \begin{aligned}
 \max_{\|\bm{\delta}\|_\infty\leq\epsilon}
 \bm{g}^{\top}\bm{\delta}
 &=\epsilon\|\bm{g}\|_1,\\
 \bm{\delta}_{\mathrm{FGSM}}
 &=\epsilon\operatorname{sign}(\bm{g}).
 \end{aligned}
 \label{eq:security-fgsm}
\end{equation}
每个坐标分别选取与梯度同号的端点，就达到线性化目标的最大值；实际输入还必须满足像素范围等合法性约束。FGSM的贡献是把输入优化化为一次反向传播，而非证明这一扰动就是原非线性损失的全局最优解\scite{goodfellow2015adversarial}

对于输入扰动，若允许集合为$\Delta$，可以搜索
$\max_{\bm{\delta}\in\Delta}\ell(f(\bm{x}+\bm{\delta}),y)$。
在可微条件下，投影梯度搜索反复沿增加损失的方向更新，再投影回$\Delta$。它是寻找失败的工具；有限步搜索失败不能认证不存在更好的扰动。

与FGSM只在原始输入附近作一次近似不同，多步搜索会在更新后的输入处重新计算梯度。记$\bm{g}^{(t)}=\nabla_{\bm{x}^{(t)}}\ell(f(\bm{x}^{(t)}),y)$，对于允许的输入集合$\bm X(\bm{x})$，无穷范数攻击常用
\begin{equation}
 \bm{x}^{(t+1)}
 =\Pi_{\bm X(\bm{x})}
 \left[
 \bm{x}^{(t)}
 +\eta\operatorname{sign}(\bm{g}^{(t)})
 \right],
 \label{eq:security-pgd}
\end{equation}
其中$\Pi$是投影，$\eta$是步长，初始化可以在允许区域内随机选取。这一投影梯度攻击通常称为PGD攻击；PGD沿用了投影梯度下降（projected gradient descent）的缩写，此处则沿增大损失的方向搜索。随机重启有助于探索不同区域；增加迭代次数则细化同一条搜索路径。二者不能相互替代，也不能因某个设置没有继续改善，就认定已经找到最坏输入。

\paragraph{针对防御机制调整搜索}

测试时应检查不同初始化、搜索策略和必要的迁移攻击。若防御包含随机性或不可微处理，攻击也需要与这些机制匹配。仅让原有梯度攻击失效可能产生\term{梯度掩蔽}（gradient masking），即梯度信号变得难以利用，却没有消除真实的决策脆弱性。

搜索失败还可能来自目标函数选择不当。Carlini与Wagner在评估防御蒸馏时，用分类前的打分差构造攻击损失，避免只依赖可能饱和的输出概率。对于目标类别$k$，一个代表性形式是
\begin{equation}
 \begin{aligned}
 \min_{\bm{x}'\in[0,1]^d}\quad&
 \|\bm{x}'-\bm{x}\|_2^2+c\,h_k(\bm{x}'),\\
 h_k(\bm{x}')={}&
 \max\left\{\max_{j\ne k}
 [z_j(\bm{x}')-z_k(\bm{x}')],-\kappa\right\},
 \end{aligned}
 \label{eq:security-cw}
\end{equation}
其中$z_j$是类别打分，$c>0$协调扰动与攻击目标，$\kappa\geq0$要求目标类别具有相应优势。其评估意义在于：原攻击的数值困难不能充当防御成功的证据，换用合适的搜索目标可能重新发现脆弱性\scite{carlini2017evaluating}

不可微预处理也有类似问题。Athalye等指出，前向计算中的量化或离散操作可能破坏梯度，却保留可被攻击的决策边界。反向传播近似法在前向仍执行真实防御，在反向用可微近似提供搜索方向；随机防御则可对多次随机变换的损失求平均，估计期望梯度。最终候选输入必须回到真实防御上检验，而不能把近似模型中的成功当作实际成功\scite{athalye2018obfuscated}

这些搜索改动处在不同层次。FGSM与PGD主要区别在更新一次还是反复更新；CW重新选择要优化的目标与扰动代价；反向近似和随机变换平均解决的是如何为含防御的系统取得可用搜索方向。它们不是互斥的攻击类别：同一测试可以同时使用多步更新、打分差目标和随机变换平均。应根据模型接口和防御机制组合这些选择，并把最终候选送回真实系统判定。

\paragraph{互补搜索与实际输入条件}

单个攻击也可能对某种损失或模型结构不敏感。AutoAttack将不同目标和搜索机制组合起来，包括自动调整步长的梯度攻击、寻找边界的攻击，以及不依赖梯度的查询攻击，用互补性减少手工调参和单一搜索造成的漏检。所谓“免参数调节”指其评估流程预设了相应规则，并不表示测试没有预算，也不表示它对任意防御都自动适用\scite{croce2020autoattack}

在一个构造的测试例子中，若两个攻击分别使100个样本中的12个和15个失败，其中8个重合，联合测试发现的是19个失败，而不是二者平均的13.5个。对于固定的一组攻击，每个样本都必须通过全部搜索才算未被攻破。该结果仍只是“这些攻击没有找到反例”的比例；真正的最坏情况鲁棒准确率不会高于这个比例。

数字空间中的成功也不等于物理环境中的成功。例如，攻击者提交的图案还可能经过拍摄、缩放和压缩，这些过程会改变输入。测试若只在直接送入模型的数组上进行，其结论就只适用于该入口。需要研究物理攻击时，应把成像距离、角度和光照等条件纳入评估；若把它们建模为随机变换，还须说明使用的变换分布为什么能够代表实际条件。

报告结果时应固定扰动集合与预算，在原本预测正确的样本中计算攻击成功率，并按扰动强度展开。对于定向攻击，成功必须是指定类别被输出；对于非定向攻击，才是任意错误均计入。若同一原样本接受多种搜索，按该样本是否至少被一种搜索攻破汇总，而非把各次查询当成独立测试样本。

\subsection{训练过程的评估}
\label{sec:security-training-test}

输入攻击只需运行被测模型；训练投毒则必须运行受到污染的学习过程，检查改变训练材料后究竟学到了什么。这里需要同时描述攻击者如何选择污染内容，以及学习者采用什么训练配置。

一次受控试验保留两份相同的训练流程：一份使用原数据，另一份只在威胁模型允许的位置加入或修改污染数据。初始化与超参数采用配对设置，并在多个随机种子下重复；留出评价集既不参与污染构造，也不用于选取最有利的防御结果。攻击者若只知道替代模型，就先在替代训练流程中选择污染内容，再检验它能否迁移到真实训练流程，不能把真实模型梯度暗中加入其知识条件。

若攻击者能改变训练集，优化对象就多了一层。设$\widetilde D$是允许提交的污染数据，
$D$是正常训练集，$\mathcal P$规定污染预算。下面先固定一种白盒训练设定：攻击者知道
$D$、真实训练算法、目标损失和随机种子分布，但在提交$\widetilde D$时还没有观察到
本次训练实际使用的种子。用
$\mathcal T(D\cup\widetilde D;\xi)$表示学习者固定的训练算法，其中$\xi\sim Q$包含
按预先规定分布抽取的初始化、数据顺序和其他训练随机性，则在攻击者选择污染数据时
尚未观察$\xi$的设定下，定向投毒可以写为
\begin{equation}
 \begin{aligned}
 \max_{\widetilde D\in\mathcal P}\quad&
 \mathbb E_{\xi\sim Q}
 \left[L_{\mathrm{target}}\bigl(\bm{\theta}_{\mathcal T}
       (\widetilde D,\xi)\bigr)\right],\\
 \text{其中}\quad&
 \bm{\theta}_{\mathcal T}(\widetilde D,\xi)
 =\mathcal T(D\cup\widetilde D;\xi).
 \end{aligned}
 \label{eq:security-poison-bilevel}
\end{equation}
这里$L_{\mathrm{target}}$按攻击目标定义，数值越大表示越接近攻击者希望造成的损害；
它不必等同于正常任务的训练损失。外层选择污染内容，训练算法映射描述学习者的响应。
$\mathcal T$必须固定优化器、超参数、停止规则及多解时的选择方式；确定性训练可以省略
对$\xi$的期望。若攻击者能够观察随机种子，或评价目标采用最坏种子、分位数而非期望，
就需要相应改变信息条件和外层统计量。实际实验应据此报告多个随机种子的分布，而不能
把多值$\argmin$直接当成一个确定响应。若攻击者只知道替代训练流程，则应像前段所述，
先在替代流程上构造污染数据，再用真实的$\mathcal T$评价迁移效果，不能直接使用这里的
白盒目标。输入攻击改变一次预测，投毒则通过学习过程改变后续多次预测。

BadNets展示了外包训练和模型供应链中的后门风险：攻击者让带触发模式的样本指向目标类别，同时维持正常样本上的分类表现。其关键区别在于攻击行为由条件触发，普通验证集可能根本没有覆盖这一条件。后门也可能随预训练模型进入下游任务，因此仅检查下游数据来源并不足够\scite{gu2017badnets}

投毒样本的标签正确，也不意味着训练过程安全。Poison Frogs研究的\term{干净标签投毒}（clean-label poisoning）保留人工看来合理的标签，却使污染样本在模型表示中靠近指定目标样本，诱导分类边界发生不利移动。这类工作说明，标签复核与训练安全审查保护的是不同环节。其成效依赖表示、训练方式及攻击者对目标的了解，不能直接推断任意少量数据都能控制任意模型\scite{shafahi2018poison}

BadNets式后门与干净标签投毒的构造目标不同。前者把“触发模式—目标类别”的关联放入训练，使多种输入在带触发模式时出现相同偏离；后者可以只针对一个指定样本，保持污染样本的表面标签正确，通过改变表示附近的训练约束影响该目标。二者都修改训练材料，但前者的验收对象通常是一组触发输入，后者还需要观察预先指定的目标是否改变判断。采用哪一种构造，应由威胁模型中的数据权限和目标范围决定。

投毒实验还需要一个没有污染的对照训练过程。若目标样本在正常训练中就经常误判，不能把污染后的错误全部归因于攻击。应在匹配的训练配置下比较目标错误率、干净数据性能与不同随机初始化的变动，并分别报告污染比例和污染样本能否实际进入训练集。后门评估需要同时检查触发条件下的目标成功率与无触发条件下的性能；只报告二者之一，会遗漏攻击有效性或隐蔽性。

例如，预先选定目标类别$k$和一组真实标签不为$k$的留出输入，对每个输入分别保留原版本和带触发模式的版本。记录带触发版本被判为$k$的比例、无触发版本的错误率，以及未污染训练模型在相同两组输入上的结果。若污染模型与对照模型只在触发版本上出现明显差异，才有证据把目标偏离与训练污染联系起来。还应按污染预算和触发强度绘制这些指标，而不是只报告最有效的一次训练。

训练数据审查与受控投毒试验应配合使用。前者检查来源、标签和表示分布等可疑线索，后者直接测量这些材料进入学习过程后是否造成损害。若需要比较两种网络结构是否更容易受投毒，应尽量匹配训练数据、算力预算与攻击权限，并让攻击分别适应两种结构；仅比较一次正常训练的准确率无法回答这个问题。对已经交付、无法重新运行其原始训练过程的模型，则转向下面的内部诊断，并说明缺少训练记录造成的证据限制。

\subsection{模型内部机制的评估}
\label{sec:security-model-test}

从模型内部评估，是利用计算图、参数、梯度或中间表示寻找薄弱机制，或者验证明确的行为性质。这需要相应的审查权限；只有查询权限时，应报告黑盒结果及其限制，不能假定自己已经完成内部检查。审查者所用的内部权限与真实攻击者可获得的权限应分别报告。

网络层数、参数数量和架构名称本身不能判定模型安全；相同结构配上不同训练数据和参数，可能具有不同边界与后门行为。前面的输入测试观察允许输入能否造成失败，内部分析则帮助解释失败如何产生、搜索是否充分，以及在哪些条件下能够排除失败。二者可以共同处理同一输入攻击；这里单列内部评估，是为了说明它额外提供的机制和验证证据。

\paragraph{计算路径与局部敏感性的诊断}

有模型访问权限时，可以检查预处理、各层计算、随机分支和输出变换是否都进入被测路径，再观察梯度或激活变化。局部输入梯度能够提示搜索方向，但单点梯度很小可能只是饱和或梯度掩蔽；它不是整个邻域内的变化上界。比较中间表示也可以发现候选异常，不过仍需回到实际预测验证它是否对应攻击目标。

若怀疑某个组件导致失败，可以在受控副本中替换或暂时移除它，保持其他条件一致，再比较攻击与正常任务。例如，分别检查加入量化前后的真实攻击结果和梯度可用性，可以帮助区分减少脆弱性与仅阻碍搜索。这样的组件比较提供机制证据，不能把修改后模型的安全结论直接当成原模型的结论。第~\ref{chap:trust-interpretability}章的计算过程解释可用于寻找候选机制，安全评估还要补上有效攻击或性质验证。

\paragraph{异常触发机制的诊断}

若模型已经取得，是否存在一种触发模式能把多类正常输入推向同一目标？Neural Cleanse按目标类别反推低复杂度触发模式，并比较不同类别所需触发规模；异常容易到达的类别可能提示后门。随后可以检查相关样本、过滤触发输入或修复模型。它把无法穷举的触发模式搜索转化为一个可操作的诊断问题，但依赖触发结构和异常性假设。输入相关触发器、多个受影响类别，以及本来就存在的类别不对称，都可能削弱这种判据\scite{wang2019neuralcleanse}

触发模式搜索用于诊断已取得模型，不等同于证明发现了真实后门。检验候选触发器时，应在独立输入上测量目标类别出现比例，比较触发前后的行为，并核查其是否满足威胁模型中的可实现条件；怀疑来自训练污染时，还需追踪训练与模型来源。没有找到候选触发器，仅说明本次搜索和触发结构没有发现异常。相关线索如何用于来源审查与参数修复，将分别在第~\ref{sec:security-training-defense}节和第~\ref{sec:security-model-defense}节讨论。

\paragraph{利用计算结构验证行为约束}

内部信息也可以用于形式化验证：把允许输入、网络计算和不希望发生的输出写成约束，检查是否存在满足全部约束的反例。对于含ReLU的网络，线性层给出线性等式，ReLU的激活与不激活状态给出分段约束。Reluplex研究了这种网络性质验证，将神经网络的计算规则纳入约束求解，而不是只在有限候选输入上试验\scite{katz2017reluplex}

例如，令$z_j(\bm{x}')$为类别$j$的打分，$y$为原输入$\bm{x}$的正确类别，$\bm X(\bm{x})$为本次允许输入集合。可以检查是否存在$\bm{x}'\in\bm X(\bm{x})$及$j\ne y$，使$z_j(\bm{x}')\geq z_y(\bm{x}')$。若有经过核验的解，它给出一个竞争类别至少追平正确类别的候选，实际误分类还应按系统的平分处理规则确认；若可靠求解器证明该约束不可满足，就排除了集合内所有这类候选。超时或返回未知，只能保留为未判定，不能计入通过。

另一种做法是沿网络计算传播有效的上下界。若能对每个竞争类别证明整个集合内$z_y-z_j$的下界严格大于零，就能认证分类不变；若下界不大于零，可能只是界太松，仍不等于找到攻击。第~\ref{sec:robustness-certification}节讨论认证的条件与报告方式。这里使用结构是为了验证具体参数下的具体性质，证明范围不自动扩展到训练投毒、自然语言语义改写或工具执行权限。

\subsection{查询接口的评估}
\label{sec:security-information-test}

模型服务的攻击目标可能是复制功能，而非使单次输出错误。Tramèr等研究了通过预测接口提取模型，说明输出概率等辅助信息可能把查询转化为关于参数的方程。例如，已知服务是逻辑回归并返回精确概率$p(\bm{x})$时，
\begin{equation}
 \log\frac{p(\bm{x})}{1-p(\bm{x})}
 =\bm{w}^{\top}\bm{x}+b.
 \label{eq:security-model-extraction}
\end{equation}
在预处理已知、概率严格介于零与一且查询设计满秩的理想条件下，$d+1$个独立方程即可确定$d+1$个未知量。实际接口的舍入、隐藏特征和更复杂模型会改变提取成本；只返回类别也会丢失方程信息，但攻击者仍可通过查询标签训练替代模型\scite{tramer2016stealing}

在一维教学例子中，查询$x=0$得到$p=0.5$，查询$x=1$得到$p=e/(1+e)$。把它们代入式~\eqref{eq:security-model-extraction}，得到$b=0$和$w+b=1$，因而恢复$w=1$。参数恢复后还要用没有参与求解的新输入核对输出，排除接口假设不成立的情况。这个例子所需的不是“高置信查询”本身，而是精确概率经对数几率变换后给出的线性方程。

若服务只返回类别，就无法从返回值计算这些方程中的对数几率。此时一种测试过程是准备或生成一批允许提交的输入，记录服务返回的标签，以这些输入标签对训练替代分类器，再用剩余预算补充当前替代模型容易分歧的区域。直接方程求解与替代模型训练都是复制功能的手段，区别在于接口提供的信息与采用的模型假设。涉及私人训练内容的提取，则另按上一章的秘密定义与攻击判据评价。

因而，模型窃取测试需要分别测量参数恢复与功能相似性。两个神经网络可以参数不同而预测接近；只检查权重是否一致会漏掉功能复制。相似性还取决于评价分布：正常用户常见输入上的一致率、决策边界附近的一致率和罕见输入上的一致率可能相差很大。查询数量、单次输出信息量和替代模型训练成本共同构成攻击预算。

具体测试可以在规定查询预算内让攻击者训练替代模型，再用独立于查询集的任务输入测量其与原服务的一致率及任务效用，画出预算增加时的变化。若攻击目标是个人信息或训练内容，成功判据则应切换为第~\ref{sec:privacy-verification}节的成员推断或内容恢复指标，而不能用功能一致率替代泄露证据。接口上观察到的每一项收益，都应连同输出粒度、查询限制和攻击者已知信息一起报告。

服务可用性还需在接口层单独测试。在隔离环境中固定并发量、单次输入规模、超时和总资源预算，对比正常请求与可疑请求的处理开销，观察正常任务完成率、等待时间和单位任务资源消耗。若延迟来自正常流量同样会触发的容量上限，应如实报告；若少量允许请求就能引发持续异常消耗，则需要定位到相应解析、推理或调度路径。功能复制、私人信息泄露和资源耗尽共享接口，但使用不同的成功与损失指标。

\subsection{交互与执行链的评估}
\label{sec:security-interaction-test}

对于生成系统，还必须把“输出了攻击者想要的内容”和“执行了攻击者想要的行为”分开。Greshake等研究间接提示注入，展示了第三方内容可以借助检索与工具调用进入模型上下文，使原本的数据来源成为行为操纵入口。这一问题依赖信息来源和授权边界，不能仅靠检查用户直接输入是否包含恶意请求来覆盖\scite{greshake2023indirect}

自动化搜索也改变了越狱评估的强度。Zou等的贪心坐标梯度方法用词元梯度筛选可能降低攻击目标损失的替换候选，再通过实际前向计算选择离散修改；在多个请求和模型上联合优化，可以研究攻击的迁移性。这里的优化目标通常是某种不期望响应的代理目标，代理损失降低仍需由最终行为评估确认。可迁移也不等于能跨越所有模型版本和系统防护\scite{zou2023universal}

AgentDojo把评估从孤立文本推进到带工具和状态的任务环境，同时检查用户任务能否完成，以及攻击者能否诱导指定偏离。这样的设计能够区分真正维持任务能力的防御与一概拒绝工具结果的防御，也便于在同一任务上比较不同注入位置。其结果仍受环境、工具权限和攻击集合约束，不能直接替代真实部署中的授权审计\scite{debenedetti2024agentdojo}

测试单位应是一次完整任务过程：在规定位置注入内容，保留模型读取、提出动作、执行端校验和状态变化，分别核对用户任务完成情况与攻击者目标。对同一任务设立无注入对照，再比较信息来源、工具权限和攻击预算变化后的结果。脱离工具执行的文本测试可以检查模型是否受到诱导，却不能独立测量实际越权损失。

沿用归档邮件的任务，可以在测试开始时固定授权邮件集合$A$，并保存收件箱初始状态。让攻击者只能修改一封待读邮件的正文；系统读取后，记录模型提出的归档集合$B$、执行器实际接受的集合$C$和最终邮件状态。$B\not\subseteq A$表示模型提出越权请求，$C\not\subseteq A$表示执行端允许越权，两者分别计数。合法任务是否完成，则根据应归档集合是否被正确处理判定。把这三项放在一起，便能区分模型未受诱导、模型受诱导但执行器拦截，以及系统通过全部拒绝维持低攻击成功率这三种情况。

评判器本身也可能误判。关键字匹配会把引用攻击内容与执行攻击混在一起，模型裁判则可能受待评文本影响。对于工具系统，优先检查外部可验证结果，例如是否向未经授权的接收者发送信息、是否修改了越权对象。无法自动判断的案例需要独立复核，并报告判定标准与不确定性，而不是用一个含义含混的安全分数掩盖。

\begin{example}[工具权限改变攻击后果]
\label{ex:security-permission}
在一个构造的邮件助手中，同一注入文本让模型提出“归档全部邮件”。系统甲允许模型直接调用无范围限制的归档接口；系统乙只向本次任务暴露经过授权的邮件标识，并在执行端复查。两者可能生成同样的文本，但只有甲能够将该文本转化为越权行动。评估必须区分模型受诱导与系统造成损害。
\end{example}

执行链检查可以把模型与执行器暂时分开。在隔离环境中直接提交正常请求、混入未授权对象的请求和超过数量上限的请求，核对底层状态，再接入受注入影响的模型重测。前一组检查规则本身，后一组检查组合后是否存在旁路。来源标记、工具路由、持久记忆和重试机制也应纳入路径清单，避免只审查一轮模型回答。

\section{攻击防御与损害控制}
\label{sec:security-defense}

评估找出损害路径后，防御要明确改变路径上的哪个环节。本节按干预位置分类：数据与模型来源保护限制有问题的材料进入系统，训练与修复改变模型如何学习或使用这些材料，输入与输出保护约束运行时的信息和资源，执行权限控制限制预测建议能够引发的动作；监测、止损与恢复处理剩余失败。每类措施都有对应的测试对象，也可能与其他措施共同阻断一条攻击链。

例如，一封含恶意指令的邮件可以在输入检查时被识别，也可能被模型正确忽略；即使这两层都失效，独立执行器仍可拒绝越权归档。若邮件已被处理并写入持久状态，则还需要隔离和恢复。分类依据是防护作用的位置与对象，而不是算法名称，不能把“用了对抗训练”作为数据来源、查询接口和工具权限都安全的证据。

\subsection{数据与模型来源的保护}
\label{sec:security-training-defense}

数据与模型来源保护决定什么材料可以进入学习或部署。需要检查数据来源、提交主体、标注与变更记录，以及模型文件和配套处理代码的版本。在签名密钥与主体身份的关联可信、验证通过的条件下，数字签名可以支持来源与完整性核验；文件摘要则需要与可信渠道提供的预期值比较，才能检查内容是否发生变化。单独提供一个摘要不能认证发布者身份，来源和完整性核验也不能证明内容没有投毒或后门。

训练准入可以实现为一条独立于模型预测的检查链：记录数据与模型来源，核对允许提交的主体和数量，将新来源暂存审查，再把获准对象纳入可追踪的训练清单。测试这条链时，应分别注入未授权来源、超过配额的提交，以及来源合法但内容受污染的材料。前两类检查权限规则是否执行，后一类才检验内容审查是否发现污染；应同时统计合法材料被拒绝的比例。

内容检查还需寻找来源合法但行为异常的材料。标签抽查可以发现部分错误标注，重复与来源分布检查可以发现异常集中提交，模型表示中的异常分析则试图发现表面标签正常的污染。不同检查对应不同线索，都要与受控污染试验和合法样本误拒情况配合评价。

谱特征方法（Spectral Signatures）先用待审数据训练候选模型，再分析同一类别中样本的内部表示：对表示中心化，找到变化最大的方向，再按样本在该方向上的投影大小筛查候选污染，清理后重新训练。Tran等的工作利用了所研究后门在表示空间留下的可分离结构；该依据依赖污染比例和表示分离条件，并不是任何后门都会显现同一种异常\scite{tran2018spectral}

这类筛查需要模型参与，可以把候选训练与清理放在隔离环境中，筛查完成后再决定哪些数据进入后续训练。它属于借助模型表示改进数据质量的循环，不是仅凭入库字段就能完成的检查；候选模型也不能因被用于审查就直接进入生产服务。

第~\ref{sec:security-model-test}节的后门诊断可以帮助选择需要审查的模型、样本和触发模式。诊断本身不会清除后门；修复需要明确改变了哪些训练材料、参数或输入处理规则，再重新测量触发条件下的攻击成功率。

异常检测结果应作为调查线索，而不是直接删除一批数据的充分理由。稀有但合法的群体也可能在表示空间中显得异常，这与第~\ref{chap:trust-fairness}章讨论的数据代表性问题相连。修复后既要验证触发行为是否消失，也要检查正常类别和稀有群体的性能，避免把防御代价集中到某些用户。

\subsection{模型训练与修复}
\label{sec:security-prediction-defense}
\label{sec:security-model-defense}

模型层面的防御改变学习规则或已学到的参数。这里按需要抑制的影响分为三种做法：通过对抗训练降低输入操纵的收益，通过聚合规则限制恶意训练更新，通过修复已交付模型削弱被诊断出的异常机制。它们可以组合，但输入扰动训练并不自动解决训练投毒，后门修复也不能代替训练来源审查。

\paragraph{针对输入操纵的训练}

对输入扰动，对抗训练把搜索得到的不利样本纳入优化，近似求解最小化最坏损失的问题。Madry等以鲁棒优化形式连接训练目标和攻击搜索。实际保护范围仍受到允许扰动、内层搜索强度和训练分布的限制\scite{madry2018adversarial}

具体训练在每个小批次中交替完成两步：固定当前参数，用第~\ref{sec:security-prediction-test}节的搜索得到允许集合内的高损失输入；固定这些候选输入，再更新模型参数以减小相应损失。以训练分布中的输入与标签取期望，其理想目标为
\begin{equation}
 \min_{\bm\theta}\ \mathbb E_{(\bm x,y)}
 \left[\sup_{\bm\delta\in\Delta(\bm x)}
  \ell(f_{\bm\theta}(\bm x+\bm\delta),y)\right].
 \label{eq:security-adversarial-training}
\end{equation}
内层搜索代表攻击者怎样利用当前模型，外层更新代表防御怎样改变这些薄弱点。实际实现通常以有限步搜索近似内层，再将得到的扰动作为当前更新的训练输入；因此模型改变后，需要重新搜索，不能永久复用同一批攻击样本。

这与把固定噪声加入训练集有实质区别：内层搜索随当前模型变化，外层训练不断响应当前薄弱点。但内层搜索太弱时，模型可能只学会抵抗训练中出现的搜索模式。因此，训练攻击与最终评估攻击不应完全相同；还要检查干净样本性能、不同扰动预算及未参与调参的攻击。

第~\ref{sec:robustness-local-enhancement}节的TRADES目标（式~\eqref{eq:robustness-trades}）进一步明确了正常预测与局部一致性之间的权衡，随机平滑的半径（式~\eqref{eq:robustness-smoothing-radius}）则在条件成立时提供认证。两类方法在这里分别承担经验增强与限定集合内保证的角色。像素范数下的保证不会自动覆盖邮件的语义改写，输入扰动下的训练也不会自动阻止训练投毒；转换攻击面时必须重新检查假设。

\paragraph{针对恶意更新的训练规则}

协作训练中，即使参与者身份合法，其提交的更新也可能被操纵。鲁棒聚合用对极端值较不敏感的规则替代直接平均；Yin等分析了基于逐坐标中位数和截尾均值的分布式学习，给出相应统计与优化条件下的误差结论。这类保证需要匹配恶意参与者比例及正常更新的分布假设\scite{yin2018byzantine}

对于参与者提交的向量更新，鲁棒聚合改变的是合并规则。例如逐坐标中位数在每一维取参与者值的中位数，而不是算术平均。一个构造的一维例子中，五个更新为$0.9,1.0,1.1,8,9$，后两个受攻击者控制；平均值为4，中位数为1.1。该例说明，中位数限制了少数极端更新对这一坐标的拉动。它不证明高维聚合后的向量一定有用：各坐标的中位数可能来自不同参与者，合法更新也可能因数据差异而分散。评价应固定参与者数量和恶意比例，比较攻击后的任务损失、正常训练收敛及不同数据群体的损失。

截尾均值则在每一坐标去掉规定数量的最大和最小值，再平均剩余更新。去掉多少必须与威胁假设对应；攻击者若能伪造大量身份，或正常客户端数据差异很大，原来的比例与集中性前提可能失效。第~\ref{sec:privacy-collaborative-computation}节的安全聚合隐藏个体更新，而这里的鲁棒聚合需要对更新施加特定运算；组合二者时应明确协议是否支持这些运算，不能把两个名称并列就当成已经得到联合保证。

\paragraph{已训练模型的修复}

如果检查发现可信的后门线索，可以清理受污染训练材料并从可信状态重训，也可以尝试修改已交付模型。Fine-Pruning结合剪枝与干净数据微调：依据干净输入下的激活削减部分低活动神经元，再用干净数据调整剩余参数。Liu等研究了单独剪枝、单独微调和二者组合的效果，以及针对剪枝设计的攻击，说明修复需要考虑攻击者的适应\scite{liu2018finepruning}

修复依赖用于诊断和微调的数据是否可信、是否覆盖正常能力，以及异常行为能否与正常机制分离。若可疑机制与有用计算共用表示，剪枝可能损害正常任务，后门也可能在未测试条件下残留。验收应重新运行触发器搜索、已知攻击和保留任务测试；仅让一个已知触发输入失效，不能证明未知触发方式都已消除。

这与第~\ref{sec:privacy-unlearning}节的机器遗忘有关，但目标不同：安全修复要求规定攻击下的异常行为受到限制，数据遗忘则与移除指定记录后的训练结果比较。修改参数可以用于两者，却需要各自的验收标准。若依赖范围无法查清或修复达不到要求，应使用可信数据和可信起点重新训练，而不是继续沿用未经验证的受污染检查点。

\subsection{推理输入与查询输出的保护}
\label{sec:security-information-defense}
\label{sec:security-input-defense}

运行时防护可以作用于输入进入模型之前，也可以作用于模型回答交付之后。输入处理限制哪些内容、以什么表示进入计算；输出与查询控制限制观察者得到多少信息、能够尝试多少次。它们通常不修改训练参数，应当在包含这些处理的完整接口上重新评估。

\paragraph{输入检查与可疑请求处置}

格式、类型、长度和取值检查可以排除不合法输入，并防止输入突破事先规定的处理范围。对语义合法却具有攻击性的输入，还可以使用异常检测、已知触发模式检测或独立风险分类器，将可疑请求拒绝、转交复核或交给受限服务处理。拒绝并不表示已经正确识别攻击，报告中需要同时统计攻击漏检、正常请求误拒以及被接受请求的剩余风险。

去噪、压缩、量化或文本规范化可以改变攻击实际到达模型时的形式，但也可能删除有用信息或只使梯度难以计算。应把处理过程交给自适应攻击共同测试，使用第~\ref{sec:security-prediction-test}节的近似反向传播、随机变换平均或黑盒搜索检查是否只是梯度掩蔽。对特定输入变换观察到收益，不足以认定所有语义保持的攻击都会被消除。

对于检索和工具返回内容，应保留来源，避免让外部内容直接覆盖系统配置或获得更高权限。过滤已知恶意短语能够处理部分已知模式，角色标记能够为模型提供区分依据，但改写与未知模式仍需检验。外部信息可能同时包含正常任务所需的步骤和恶意要求，因此最终行动仍需要下一小节的独立授权检查。

\paragraph{查询输出与资源配额的控制}

查询限制和输出粒度控制可以增加模型窃取或信息提取的成本，也可能影响正常使用。应结合第~\ref{chap:trust-privacy}章判断是否需要输出层面的正式保护，而不能仅以接口隐藏了参数便宣称信息不会泄露。

输出粒度控制直接改变攻击者每次查询得到的约束。对于式~\eqref{eq:security-model-extraction}中的逻辑回归，精确概率给出等式；若只发布一个概率区间，攻击者只能知道对数几率落在相应区间；若只返回类别，则只得到决策边界某一侧的信息。限制查询数量进一步约束这些信息能够累积多少。机制的代价是，正常用户也可能失去校准概率或批量处理能力，因此需要按业务用途决定哪些信息应当返回。

可以在同一被测模型上比较精确概率、有限精度概率和仅类别三种接口，分别在相同查询预算下训练替代模型，并在独立输入上测量功能一致率；再固定正常用户的任务，检查输出变化造成的损失。这样能够看到防御究竟减少了接口信息，还是仅拖慢了所有使用者。即使只返回类别，攻击者仍可能通过较多查询近似分类边界，所以评价结果应是一条预算与提取效用的关系，而不是有无概率字段这一布尔判断。

隐私保证与模型保密性也不能互相替代。差分隐私限制单个参与者对发布结果的影响，却允许发布本身包含有价值的总体规律；攻击者仍可能近似复制这些规律。反过来，隐藏模型结构也不能证明训练数据受到保护。接口设计必须先明确要保护个人信息、模型商业价值，还是对安全决策规则的了解程度，再选择相应评价目标。

控制措施实施后，应按第~\ref{sec:security-information-test}节重新比较相同预算下的提取收益，并测量正常用户被限流、输出精度下降或响应延迟的代价。若攻击者可以使用多个账户或多次会话，预算控制还需与威胁模型中的身份边界一致；单一会话的限制并不自动约束整个攻击过程。

对资源消耗还可以限制单次输入规模、并发请求、生成长度、工具调用次数和任务总时长，并在队列中隔离不同租户。限制应作用于完整任务和可识别的主体，而非仅某次调用；否则一次任务中的重试或多会话累积仍可能超出原定预算。测试需要同时观察是否抑制异常消耗，以及长文本、批处理等正常需求是否受到不合理限制。

\subsection{行动执行中的权限控制}
\label{sec:security-execution-defense}

行为约束中的执行干预把模型输出限制为动作提案，再由独立执行器核验授权对象、参数与当前状态，具体机制见第~\ref{sec:alignment-intervention-execution}节。对抗安全进一步检查攻击者能否绕过这些条件。安全结论不能停留在“存在权限检查”，而要覆盖攻击者能够控制的身份、内容、调用顺序和反馈。

最小权限仍是测试的起点：摘要任务不应获得删除能力，读取指定文档不应获得浏览全部私人目录的权限。攻击者可能伪造对象标识、把越权对象藏入批量参数、利用检查与执行之间的状态变化，或借多个账户和连续低权限动作绕过单次上限。沙箱和结构化参数只有在这些路径都受约束时才有明确作用，名称本身不能说明隔离范围。

多层防御应检查失效是否独立。如果输入过滤、行为评判和最终授权都依赖同一个容易被相同文本操纵的模型，增加层数未必显著减少风险。防御设计需要追踪各层共享的信息与弱点。

测试时应直接向执行器提交伪造主体、越权对象、重复对象、超过数量限制的批量请求及状态已经改变的请求，再把同样条件放回受提示注入影响的完整交互。前一组检查授权规则本身，后一组检查模型、路由和执行器组合后是否存在旁路。每次测试都要核对底层状态和正常任务完成情况；错误文本被拒绝只是行为结果，只有越权对象始终没有发生状态变化，才支持执行边界在相应攻击条件下有效。

检索增强系统还应保留来源信息。来自网页的操作建议可以成为用户任务所需的事实材料，却不应因此修改工具权限。把文本放入引号或标为“不可信”，有助于训练和推理时区分内容角色，但这些仍是模型需要正确理解的信号。真正限制损害的条件是，忽略这些信号时也无法越过外部授权检查。涉及不可逆动作时，复核界面还要显示具体对象和后果，让复核者能够判断所批准的究竟是什么。

\subsection{攻击监测与损害恢复}
\label{sec:security-recovery}

事前防护仍可能被突破，运行中的检测与应急处置需要衔接。可以监测异常查询密度、资源消耗、工具调用范围、连续拒绝、模型版本变化和外部信息流向，再结合实际状态与独立记录判断是否需要响应。这些信号提供异常线索，不能只凭模型自报“没有受到攻击”确认安全；误报与漏报也应通过演练测量。

检测到高风险迹象后，可以暂停相关任务与后续工具调用、撤销会话凭据、隔离可疑数据和模型版本，并将服务切换到受限功能或经过验证的版本。不同动作应按损害传播位置选择：正在外发信息时要及时关闭出口，模型已受污染时还要停止相关版本继续产生决策。只拦住当前输入，不能清除训练参数或持久记忆中的污染。

有些攻击只能在造成部分影响后被发现。限制单次动作规模、保存必要恢复状态、及时撤销权限和停止后续动作，可以减少损失范围。恢复机制也需要测试：能够恢复文件不意味着能够收回已经发送的信息，能够停止后续动作也不能撤销此前不可逆操作。

恢复演练可以在隔离测试环境中记录初始状态，执行一条会错误修改邮件标签的构造轨迹，在规定时刻触发检测，再观察停止后续动作、撤销权限和恢复标签是否成功。测量应包括检测至停止的时间、停止前影响了多少对象、恢复了多少状态，以及仍无法恢复的损失。对已经发送到外部的信息，测试结果只能记录进一步传播是否受限，不能把本地删除日志当成收回信息。这个过程检验的是发现失效后系统还能控制什么，而不是重新评价最初的攻击是否成功。

恢复还涉及证据保全。系统需要记录足够的信息来区分模型错误、外部注入与授权配置问题，但日志本身可能包含敏感输入，应按第~\ref{chap:trust-privacy}章的原则控制访问与保留范围。发现攻击后若只替换模型而不关闭同一注入入口，风险可能继续存在；若只删除恶意文档而不检查已经写入的记忆和派生状态，也可能留下后续触发条件。恢复范围应沿实际信息传播路径确定。

防御完成后应回到第~\ref{sec:security-evaluation}节复测正常任务、已知攻击和针对新防御的攻击。结论应写为在指定能力与预算下哪些损害受到限制，并说明尚未覆盖的条件。

止损与恢复后的重新开放也需要判据。对可疑版本或入口，应确认处置范围、修复验证结果和仍未覆盖的攻击条件，再逐步恢复权限并继续观察。防御的共同验收条件是：在相同威胁模型下减少了哪些损害，正常任务付出了什么代价，以及失效后能在多长时间内停止传播；不是部署的检测器或过滤层数越多就越安全。

\section{前沿研究}
\label{sec:security-frontiers}

智能体的风险来自动态交互。外部内容既可能包含有用操作信息，也可能包含试图扩展权限的指令，简单地忽略全部外部指令会降低任务能力。评估因此需要同时考察正常信息使用和恶意操纵，而不能只在孤立提示上比较是否拒答。

AgentDyn围绕动态任务、有用第三方指令和更复杂的用户目标构造评测，讨论安全性与过度防御之间的取舍。这是对评估情境的扩展，论文中的结果仍受所测任务、攻击与防御集合限制\scite{li2026agentdyn}

它所指向的困难可以从前面的邮件例子看出：一封可信业务邮件可能包含完成任务必需的步骤，但其中也可能夹带超出授权的操作。完全忽略步骤会使任务失败，全部照做又会扩大攻击面。研究重点因而从“文本中有没有命令”转向“哪一部分内容在何种来源与权限条件下可以影响哪个动作”。评估需要保留这种混合情境，而非让正常任务与恶意文本在表面形式上容易区分。

另一个变化是攻击者能够长期适应防御。静态测试集容易逐渐成为训练材料，使分数改善与防护能力改善难以区分。具有明确预算的自动化红队测试可以持续提出新候选，但候选生成器本身也有搜索偏好。可信证据需要结合多种攻击来源、保留未用于防御调优的任务，并记录攻击者获得了哪些反馈；仅增加生成次数不保证覆盖新的失败机制。

多模态内容、长期记忆和工具组合进一步扩大攻击面。某段内容可能在读入时无害，却在后续被重新当作指令；多个低权限操作也可能组合出较大损害。未来研究需要从单次成功率扩展到行为链、状态变化和可恢复性，并让攻击者能够针对完整协议调整策略。

\section{本章小结}
\label{sec:security-summary}

对抗攻击通过有目的的输入选择、训练操纵、信息获取或行为诱导，使机器学习系统偏离任务与安全要求。攻击类型说明采取什么操作，完整性、机密性和可用性说明造成什么损害，威胁模型则限定攻击者能看见、改变和尝试多少。模型出错、提出越权动作与系统真正造成损失，需要分别记录。

评估中，输入搜索寻找具体反例，受控训练比较检验污染如何改变模型，内部分析定位机制或验证规定性质，查询试验测量信息与资源暴露，完整交互测试核对权限和状态变化。这些证据共同说明系统在哪些条件下脆弱；架构名称、单次攻击失败或一张干净测试成绩单都不足以替代它们。攻击成功率还应与预算、损害、正常任务代价及证据覆盖范围一起解释。

防御沿同一信息和行动路径布置。来源审查不能消除已进入参数的后门，模型训练不能代替查询配额与执行授权，运行时过滤也不能自动修复持久污染。选择措施时，应根据已发现的薄弱位置组合数据、模型、接口与执行端保护，并让攻击者针对新防御重新搜索。监测和恢复进一步决定突破发生后还能控制多少损失；信息已经外泄与可撤销的状态修改，具有不同的恢复边界。

目标对齐与行为约束还要处理没有外部攻击时出现的行为偏离，相关讨论见第~\ref{chap:trust-alignment}章。本章进一步加入有目的的操纵：攻击者可能利用预测误差、环境变化、信息泄露或权限缺口，使既有保护失效。因此，安全验证需要把前面各章的性质放回同一个系统中，检查它们在攻击者能力范围内是否仍然成立；一项性质的改进不能代替对其余损害路径的检查。

\paragraph{延伸阅读}

Madry等人的鲁棒优化与Athalye等人的自适应攻击研究，适合继续理解防御训练和攻击评估为何必须配套。前者把允许扰动写入训练目标，后者说明梯度不可用可能只是搜索受阻；阅读时应区分经验搜索失败与形式化认证\scite{madry2018adversarial,athalye2018obfuscated}

对于具有检索和工具权限的系统，可继续阅读间接提示注入与AgentDojo。前者追踪第三方内容怎样进入行为链，后者同时评价用户任务和攻击目标；两者都要求把模型输出、执行器接受的动作和最终状态分开记录\scite{greshake2023indirect,debenedetti2024agentdojo}
